温度 \(T>0\) 恒定的理想蒸气处于过饱和状态，实际分压为 \(p=Sp_{\rm sat}(T)\)，饱和比 \(S>1\)。液态分子体积 \(v_\ell>0\) 近似常数，平面液汽界面张力 \(\gamma>0\)，分子热能用 \(k_BT\) 表示。假定液滴为半径 \(r>0\) 的球体、蒸气是大热库，忽略液态压缩和液滴平移/混合熵的次级修正；表面能在经典成核近似下取 \(4\pi r^2\gamma\)。
(1) 从两相化学势平衡、理想蒸气化学势与 Laplace 压差出发，推导半径 \(r\) 的 Kelvin 饱和蒸气压关系；解释平面极限及为什么半径越小越难以稳定。
(2) 在恒温、恒外部过饱和度下，以 \(n=4\pi r^3/(3v_\ell)\) 为液滴分子数，构造成核的可逆自由能 \(\Delta G(r)\)。求临界半径、临界分子数、势垒高度，证明临界核是自由能的极大值而非稳定液滴。
(3) 证明临界核满足的 Kelvin 半径条件恰好与第 (2) 问一致。讨论 \(S\downarrow1\) 及强过饱和两极限、成核率 \(J\propto e^{-\Delta G_*/(k_BT)}\) 的敏感性，并推导势垒与临界分子数的精确比值关系。
(4) 为检验平面表面张力假设的首阶误差，改用 \(\gamma(r)=\gamma[1-2\delta/r]\)，其中 Tolman 长度 \(\delta\) 为有符号常数且 \(|\delta|\ll r_*\)。仍按给定模型把界面项定义为 \(4\pi r^2\gamma(r)\)，不另行改变体积项。推导临界半径与势垒到 \(\delta\) 一阶的修正，并判断 \(\delta>0\) 对成核指数的影响和近似的有效范围。